\subsection{Analisi delle proprietà di sicurezza}


\paragraph{Autenticità}:
Ogni Utente intenzionato a votare possiede un certificato digitale valido firmato dalla Certificate Authority. L'Authenticator verifica la validità del certificato e infine verifica la firma sul voto cifrato. Solo l'Utente in possesso della chiave privata $sk_U$ può produrre una firma valida, quindi un attaccante non può impersonare un elettore senza compromettere la sua chiave privata. \\
La proprietà è garantita solo se assumiamo che, in fase di richiesta certificato utente (2.2), tutti gli Utenti inseriscano nel campo Subject del certificato non firmato la propria matricola corretta. Se così non fosse, chiunque potrebbe fingersi qualcun'altro.

\paragraph{Unicità}:
La proprietà di unicità è garantita dall'Authenticator attraverso la struttura dati VoterMap, che mantiene traccia degli elettori che hanno già votato. In particolare, l'identificativo dell'elettore viene rappresentato mediante l'hash della sua chiave pubblica $H(pk_U)$. Quando un nuovo voto viene ricevuto, l'Authenticator verifica se tale identificativo è già presente nella mappa. In caso affermativo, la richiesta viene immediatamente rifiutata; in caso contrario il voto viene accettato e l'identificativo viene inserito nella struttura. Questo meccanismo impedisce a un utente legittimo di votare più volte utilizzando lo stesso certificato e la stessa coppia di chiavi. \\
La proprietà è garantita assumendo che l'Authenticator operi correttamente, che la VoterMap non venga alterata da soggetti esterni e che un elettore non possa ottenere più certificati validi associati alla stessa votazione.

\paragraph{Integrità}:
La proprietà di integrità è garantita mediante l'utilizzo delle firme digitali su tutti i messaggi critici del protocollo. In particolare, il voto trasmesso dall'Utente all'Authenticator è accompagnato da una firma digitale generata mediante la chiave privata dell'Utente. L'Authenticator verifica tale firma utilizzando la corrispondente chiave pubblica ottenuta dal certificato dell'Utente. Analogamente, il messaggio inoltrato dall'Authenticator al Server viene firmato dall'Authenticator e verificato dal Server. \\
La proprietà è garantita assumendo che le chiavi private degli attori onesti (Utente, Authenticator e Server) rimangano segrete. \\
L'integrità del contenuto del voto durante la trasmissione è garantita dal protocollo. Tuttavia, il protocollo assume che il Server, una volta decifrato il voto, si comporti in maniera semi-honest e non alteri volontariamente il risultato dello scrutinio. Pertanto l'integrità del risultato finale dipende parzialmente da tale assunzione di fiducia.


\paragraph{Segretezza}:
La proprietà di segretezza viene garantita mediante l'utilizzo della cifratura asimmetrica sul voto trasmesso dall'utente. Infatti l'Utente, oltre a trasmettere il voto cifrato con la sua chiave pubblica da salvare sull'Authenticator, trasmette il voto cifrato utilizzando la chiave pubblica del Server e quindi non decifrabile dall'Authenticator. L'Authenticator quindi sa che l'Utente ha votato ma non può ottenere l'informazione relativa al voto. D'altro canto, l'Authenticator inoltra il voto cifrato al Server il quale riceve esclusivamente il voto cifrato e le informazioni necessarie a verificarne la provenienza dall'Authenticator. Non riceve certificati dell'Utente, chiavi pubbliche dell'Utente o altri identificativi che consentano di risalire direttamente all'elettore che ha espresso il voto. \\
La proprietà è quindi garantita sotto l'assunzione che l'Authenticator e il Server si comportino come entità semi-honest e non colludano. Se dovessero colludere, l'Authenticator potrebbe inviare al Server anche l'identità dell'Utente votante oltre al suo voto. Questo permetterebbe al Server di associare l'identità dell'Utente al suo voto espresso.

\paragraph{Verificabilità Individuale}:
La proprietà di verificabilità individuale è garantita dalla struttura VoterMap mantenuta dall'Authenticator e dalla procedura di verifica del voto. Infatti, per ogni elettore che ha votato, l'Authenticator conserva una coppia costituita dall'identificativo dell'utente $H(pk_U)$ e dal voto cifrato con la chiave pubblica dell'utente $c = Enc_{pk_U}(v)$. Successivamente, l'Utente può autenticarsi presso l'Authenticator e richiedere il recupero del proprio voto cifrato. \\
Una volta ricevuto il voto cifrato associato al proprio identificativo, l'Utente può decifrarlo utilizzando la propria chiave privata $sk_U$ ed ottenere il voto in chiaro. Confrontando il voto ottenuto con quello che intendeva esprimere, l'Utente può verificare che il voto registrato dal sistema corrisponda effettivamente alla propria scelta elettorale. \\
La proprietà è garantita assumendo che l'Authenticator conservi correttamente la VoterMap e restituisca all'Utente il voto effettivamente associato al suo identificativo. Inoltre, si assume che il Server operi in modo semi-honest e non alteri i voti ricevuti dall'Authenticator prima dello scrutinio. In caso contrario, l'Utente potrebbe verificare correttamente il voto conservato dall'Authenticator ma non quello effettivamente conteggiato dal Server.

\paragraph{Verificabilità Universale}:
La proprietà di verificabilità universale non è completamente garantita dal protocollo proposto. Al termine dello scrutinio il Server pubblica il risultato finale della votazione e il numero totale di voti validi, firmando digitalmente tali informazioni. Chiunque può quindi verificare l'autenticità e l'integrità del risultato pubblicato mediante la verifica della firma digitale del Server. \\
Tuttavia, il protocollo non prevede la pubblicazione dei voti registrati né meccanismi crittografici che consentano a soggetti esterni di ricostruire autonomamente il conteggio effettuato dal Server. Di conseguenza, un osservatore esterno può verificare che il risultato pubblicato provenga effettivamente dal Server, ma non può verificare indipendentemente che ciascun voto sia stato conteggiato correttamente durante lo scrutinio. \\
La proprietà è quindi soddisfatta solo parzialmente e dipende dall'assunzione che il Server si comporti come un'entità semi-honest. Un Server malevolo potrebbe infatti alterare il conteggio finale dei voti senza che un osservatore esterno possa rilevare tale comportamento utilizzando esclusivamente le informazioni pubblicate dal protocollo.


\subsection{Analisi degli attacchi del Threat Model}


\paragraph{Identity Theft}:

\paragraph{Spoofing Authenticator}:

\paragraph{Spoofing Server}:

\paragraph{Man-in-the-Middle}:

\paragraph{Replay Attack}:

\paragraph{Eavesdropping}:

\paragraph{Traffic Analysis}:

\paragraph{Authenticator Curioso}:

\paragraph{Server Curioso}:


\subsection{Analisi delle assunzioni di fiducia}

